\documentclass[11pt]{amsart}
\usepackage[margin=1.0in]{geometry}
\usepackage{amsmath,amssymb,amsthm,mathtools,booktabs,enumitem,xcolor,array}
\usepackage[colorlinks=true,linkcolor=blue!50!black,citecolor=blue!50!black,urlcolor=blue!50!black]{hyperref}
\newtheorem{theorem}{Theorem}[section]
\newtheorem{proposition}[theorem]{Proposition}
\newtheorem{lemma}[theorem]{Lemma}
\newtheorem{corollary}[theorem]{Corollary}
\theoremstyle{definition}
\newtheorem{definition}[theorem]{Definition}
\theoremstyle{remark}
\newtheorem{remark}[theorem]{Remark}
\newcommand{\R}{\mathbb{R}}
\newcommand{\E}{\mathbb{E}}
\newcommand{\T}{^{\mathsf T}}
\newcommand{\cov}{\operatorname{Cov}}
\newcommand{\relu}{\operatorname{relu}}
\newcommand{\sym}{\operatorname{sym}}
\newcommand{\wick}[1]{{:}#1{:}}
\newcommand{\Ybar}{\overline{Y}}

\title[The final layer as a measurement]{The final layer as a random-feature measurement:\\ exact RMS as a kernel distance, transverse analyticity,\\ and what a real depth-16 network demands}
\author{Stage 23 of the quantum-cut / noncommutative-geometry programme}
\date{10 October 2026}

\begin{document}
\begin{abstract}
Treat the realised final layer as what it is: $n$ independent Gaussian random-feature measurements of one
state, the law $\lambda$ of the penultimate activations. Three things follow.
\begin{enumerate}[label=(\arabic*),leftmargin=1.6em,itemsep=1pt,topsep=1pt]
\item \emph{Exact RMS as a kernel distance} (Thm.~\ref{thm:mmd}). Let $\lambda'$ be any signed surrogate state
independent of $W_L$; every estimator of stages~14--22 is the readout of one. Then
$\E[\mathrm{RMS}^2\mid\mathcal F]=\mathrm{MMD}^2_{k_1}(\lambda,\lambda')$, the maximum mean discrepancy for the degree-$1$
arc-cosine kernel. The realised RMS concentrates around it at rate $n^{-1/2}$. The error ledger of stage~18 is
exactly the Mercer expansion of this distance.
\item \emph{Transverse analyticity} (Thm.~\ref{thm:analytic}). The arc-cosine kernel is singular only on
coincident directions. Expanded around the collective angle, its Taylor series in the transverse overlap
$\rho_\perp$ of two inputs has nonnegative coefficients. The tail beyond order $J$ is bounded by $\pi\,|\rho_\perp|^{J+1}$.
So for the RMS, the kink is averaged over the random final weights, and the obstruction of stage~20, which
concerns a single neuron's moments, does not apply.
\item \emph{What a real network demands} (\S\ref{sec:real}). We measured one realised network ($n=512,1024$;
$L=16$). The transverse field is strongly anisotropic: its effective dimension is $n_{\rm eff}\approx80$ at $n=1024$,
so $|\rho_\perp|\approx0.11$, not $n^{-1/2}=0.03$. The per-neuron standardised skewness is $\approx0.12$ RMS. The absolute
$\kappa_3$ is $O(1/n)$, as before, but the collapsed transverse variance ($\sigma_k^3\approx0.06$) amplifies it. This
corrects earlier size claims. The $\kappa_3$ term moves each mean by $\approx10^{-3}$ RMS, sixteen times the target,
so the history readout must be accurate to about $6\%$. The $10\kappa_3^2$ term of $m_6$, which the audit flagged,
and the incoherent $\kappa_4$ are each at the target scale and must be kept.
\end{enumerate}
We also prove that each random layer adds $1/n$ to the inverse participation ratio
(Prop.~\ref{prop:pr}). Measured, $n/n_{\rm eff}$ grows linearly in depth with slope $0.85$. So the per-neuron expansion
parameter is the depth-to-width ratio $\sqrt{0.85L/n}\approx0.12$, not $n^{-1/2}$. Two adaptive checks follow from the
probe block: an order selector from $n_{\rm eff}$, and a skewness calibration test with a slope standard error of
about $0.02$ (\S\ref{sec:adapt}).
\end{abstract}
\maketitle

\section{The final layer measures a state}
Conditional on $\mathcal F=\sigma(W_1,\dots,W_{L-1})$, the columns $w_k\sim N(0,s^2I)$ are i.i.d. and independent of the penultimate
law $\lambda$. Put $y(w)=\int(w\cdot u)_+\,d\lambda(u)$, so $y_k=y(w_k)$. For a finite signed measure $\lambda'$ independent of
$W_L$, put $\hat y(w)=\int(w\cdot u)_+\,d\lambda'(u)$.

\begin{theorem}[Exact RMS as an arc-cosine MMD]\label{thm:mmd}
Let $k_1(u,u')=\E_w[(w\cdot u)_+(w\cdot u')_+]=\frac{s^2}{2\pi}\|u\|\|u'\|\,\kappa(\cos\theta)$, where
$\kappa(t)=\sqrt{1-t^2}+(\pi-\arccos t)\,t$. Then
\[
\E\big[\mathrm{RMS}(\hat y-y)^2\,\big|\,\mathcal F\big]=\mathrm{MMD}^2_{k_1}(\lambda,\lambda')=\iint k_1\,d(\lambda-\lambda')^{\otimes2},
\]
and $\mathrm{Var}(\mathrm{RMS}^2\mid\mathcal F)=\mathrm{Var}_w\big((\hat y-y)^2\big)/n$. Moreover:
\begin{enumerate}[label=(\alph*),leftmargin=1.6em,itemsep=1pt,topsep=1pt]
\item \emph{Mercer equals ledger.}
\[
\mathrm{MMD}^2=s^2\sum_{r\ge0}\frac{c_r^2}{r!}\,\big\|M_r(\lambda)-M_r(\lambda')\big\|_F^2,\qquad M_r(\nu)=\int\|u\|^{1-r}u^{\otimes r}d\nu ,
\]
with $c_r$ the Hermite coefficients of $\relu$. This is stage~18's Wiener-chaos ledger in the final weights.
\item \emph{All our estimators are signed-state readouts.} The Gaussian closure is the readout of $N(\mu,C)$.
The cumulant corrections are readouts of the multivariate Edgeworth signed measure, whose $\kappa_r$ tensors
are the hidden cumulants. The collective form is the readout of a mixture of these.
\end{enumerate}
\end{theorem}
\begin{proof}
$\hat y-y=\int(w\cdot u)_+\,d(\lambda'-\lambda)$. Squaring, taking $\E_w$ and using Fubini gives the MMD. The $n$ columns are
i.i.d., which gives the variance. The arc-cosine formula is the classical Gaussian computation.

For (a), the Mehler expansion $(w\cdot u)_+=s\|u\|\sum_rc_r\mathrm{He}_r(w\cdot\hat u/s)/r!$ and Hermite orthogonality give
$k_1=s^2\|u\|\|u'\|\sum_r(c_r^2/r!)(\hat u\cdot\hat u')^r$. Integrate against $(\lambda-\lambda')^{\otimes2}$.

For (b), the one-dimensional marginal of a multivariate Edgeworth measure along $w$ is the one-dimensional
Edgeworth measure with cumulants $\langle\kappa_r,w^{\otimes r}\rangle$. Its $\relu$ readout is stage~17's $F$.
\end{proof}

In NCG terms, the kernel $k_1$ is the GNS inner product of the ``final-layer state''. The mean embedding
$\xi_\lambda=\int\phi_u\,d\lambda(u)$, with $\phi_u=(w\mapsto(w\cdot u)_+)$, is a vector in $L^2(\gamma_s)$. The final means are its overlaps
with $n$ random Gaussian measurement vectors. As with classical shadows, random measurements estimate a
distance between states: the RMS is the GNS distance $\|\xi_\lambda-\xi_{\lambda'}\|$ up to $n^{-1/2}$ fluctuations. The
design question becomes: find a computable signed state $\lambda'$ at the penultimate layer within GNS
distance $6\times10^{-5}$ of the true one.

\section{Transverse analyticity: why the RMS escapes the moment obstruction}
Write each penultimate vector as $u=R(\cos\beta\,\hat\mu+\sin\beta\,\hat v)$, with $\hat v\perp\hat\mu$ (stage~18), and for a pair put
$\rho_\perp=\hat v\cdot\hat v'$. Then
\[
\cos\theta=t_0+h,\qquad t_0=\cos\beta\cos\beta',\qquad h=\sin\beta\sin\beta'\,\rho_\perp .
\]

\begin{lemma}[Positivity]\label{lem:pos}
$\kappa(t)=\sum_{m\ge0}b_mt^m$ with $b_m=2\pi c_m^2/m!\ge0$ and $\kappa(1)=\pi$. For $t_0\in[0,1)$, every Taylor coefficient
$\kappa^{(j)}(t_0)/j!=\sum_mb_m\binom mjt_0^{m-j}$ is nonnegative, and $\sum_j\kappa^{(j)}(t_0)(1-t_0)^j/j!=\kappa(1)$.
\end{lemma}
\begin{proof}
The expansion is Mehler's, as in Theorem~\ref{thm:mmd}(a), normalised so that $k_1(u,u)=s^2\|u\|^2/2$. The series has
nonnegative coefficients and converges at $t=1$, so Taylor re-expansion at $t_0$ is valid on $|h|\le1-t_0$.
\end{proof}

\begin{theorem}[Transverse expansion with geometric tail]\label{thm:analytic}
For every pair, $|h|\le1-t_0$. Splitting $\kappa(t_0+h)=\sum_{j\le J}\kappa^{(j)}(t_0)h^j/j!+\mathcal T_J$, we have
\[
0\le\mathcal T_J\ \text{in the PD sense},\qquad|\mathcal T_J|\le\pi\,|\rho_\perp|^{J+1}.
\]
Consequently
$\mathrm{MMD}^2_{k_1}(\lambda,\lambda')=\sum_{j\le J}\sum_m\frac{s^2b_m}{2\pi}\binom mj\|\Delta_{jm}\|^2+E_{>J}$, where
\[
\Delta_{jm}=\int R^{1-m}P^{m-j}u_\perp^{\otimes j}\,d(\lambda-\lambda'),\qquad
0\le E_{>J}\le\frac{s^2}2\iint RR'|\rho_\perp|^{J+1}\,d|\lambda-\lambda'|^{\otimes2}.
\]
\end{theorem}
\begin{proof}
$1-t_0-|h|\ge1-\cos(\beta-\beta')\ge0$, so $|h|\le1-t_0$. Also $|h|/(1-t_0)\le|\rho_\perp|$, because $1-\cos\beta\cos\beta'\ge\sin\beta\sin\beta'$. By
Lemma~\ref{lem:pos} the tail is a nonnegative combination of powers of $h$, bounded by
$(|h|/(1-t_0))^{J+1}\sum_j\kappa^{(j)}(t_0)(1-t_0)^j/j!$. The low part is a finite sum of products of feature maps. Indeed
$R\cos^{m-j}\beta\sin^j\beta\,\hat v^{\otimes j}=R^{1-m}P^{m-j}u_\perp^{\otimes j}$, which gives the $\Delta_{jm}$.
\end{proof}

\noindent\textbf{What this means.}
\begin{itemize}[leftmargin=1.2em,itemsep=1pt,topsep=1pt]
\item The arc-cosine kernel is singular only where two hidden directions coincide, which has probability
zero for pairs of independent inputs. On the support of such pairs it is analytic, with expansion parameter
$|\rho_\perp|$.
\item The per-neuron obstruction of stage~20 concerns one fixed $w$, which probes the kink. The RMS is an
average over $n$ independent $w$, which smooths the kink into $k_1$.
\item The order-$j$ terms vanish exactly when $\lambda'$ matches the transverse moment tensors of order $j$
\emph{jointly with} the collective coordinates (all $m$). That is stage~18's $B_j(a)$ matching.
\item A certificate needs two things: bounds on the mismatches $\|\Delta_{jm}\|$ for $j\le J$, and a pair-overlap tail
bound. It needs no per-neuron cancellation theorem.
\end{itemize}

\noindent The expansion around $t_0$ assumes $P\ge0$, that is $\beta\le\pi/2$. This holds on the collective cap at depth; at
$n=1024$, $L=16$, $\sin^2\beta\approx0.07$ for every input sampled.

\section{What a real depth-16 network looks like}\label{sec:real}
One realised He network, with $2000$ Gaussian inputs for the penultimate layer and $20000$ for the final
preactivations. Script: \texttt{theory23/}.
\begin{center}\small
\begin{tabular}{@{}lcc@{}}\toprule
quantity at layer $L-1=15$ & $n=512$ & $n=1024$\\\midrule
$\sin^2\beta$: mean (sd) & $0.068\ (0.015)$ & $0.067\ (0.010)$\\
pair $\cos\theta$: mean (sd), max & $0.932\ (0.015)$, $0.979$ & $0.933\ (0.010)$, $0.973$\\
transverse overlap $\rho_\perp$: RMS ($n^{-1/2}$) & $0.157\ (0.044)$ & $0.112\ (0.031)$\\
$n_{\rm eff}=1/\E\rho_\perp^2$; participation ratio of $C_\perp$ & $40.8$; $38.1$ & $80.1$; $75.1$\\
$\mathrm{Var}\log R^2$ over inputs & $0.052$ & $0.029$\\\midrule
final preactivations: standardised skewness, mean (RMS) & $+0.006\ (0.203)$ & $+0.004\ (0.116)$\\
excess kurtosis: mean; incoherent sd & $0.155$; $0.061$ & $0.066$; $0.026$\\\bottomrule
\end{tabular}
\end{center}
At $n=256$ the RMS skewness is $0.319$.

\noindent\textbf{Reading.}
\begin{enumerate}[label=(\roman*),leftmargin=1.6em,itemsep=1pt,topsep=1pt]
\item \emph{The transverse field is effectively about $n/13$-dimensional at depth.} Its covariance is far from
isotropic, so the expansion parameter of Theorem~\ref{thm:analytic} is $|\rho_\perp|\approx0.11$, not $0.03$. The
participation ratio of $C_\perp$ is a Gaussian-part quantity, $(\operatorname{tr}C_\perp)^2/\|C_\perp\|_F^2$, costing $O(n^2)$ once $C_{L-1}$ is
known. Its agreement with $1/\E\rho_\perp^2$ makes it a computable order selector.
\item \emph{Skewness has random sign but is not small.} Its mean is near zero, consistent with stage~22's
$\E_w\kappa_3=0$. Its RMS is $0.12$. The absolute size is what that lemma predicted:
$\kappa_3=\mathrm{skew}\cdot\sigma^3\approx0.116\times0.064\approx7.4\times10^{-3}\approx7/n$, which implies $\|T\|_F^2\approx1.2n$. The standardised size is
large because $\sigma_k^2\approx s^2\operatorname{tr}C_\perp\approx2(1-\rho_{L-1})$ has collapsed at depth. So the factor $(1-\rho)^{-3/2}\approx40$
multiplies every standardised third cumulant.
\item \emph{The collective cocycle is confirmed.} $\mathrm{Var}\log R^2=0.029$ at $n=1024$, against stage~18's
trajectory estimate of $0.030$.
\end{enumerate}

\noindent\textbf{Corrections.} Stage~20's ``$c_3,c_4=O(1/n)$'' and stage~22's Lemma~2.2 are correct in absolute
units, as $n$-scaling statements. But the Edgeworth coefficients act on \emph{standardised} cumulants, and at
$L=16$ these carry the depth factor $(1-\rho_{L-1})^{-(r-2)/2}$. The audit's $10\kappa_3^2$ objection is therefore
numerically relevant here, though not as an $n$-scaling issue; see the ledger below.

\begin{proposition}[Inverse participation ratio grows by $1/n$ per random layer]\label{prop:pr}
For a fixed PSD $C$ and $W$ with i.i.d.\ $N(0,s^2)$ entries,
\[
\frac{\E\|W\T CW\|_F^2}{(\E\operatorname{tr}W\T CW)^2}=\Big(1+\frac1n\Big)\frac{\|C\|_F^2}{(\operatorname{tr}C)^2}+\frac1n .
\]
So, to leading order, a linear random layer acts by $1/\mathrm{PR}\mapsto1/\mathrm{PR}+1/n$. This is the participation-ratio form of free
multiplicative convolution with a Marchenko--Pastur law. Starting from $\mathrm{PR}(I)=n$, the linear part predicts
$n/\mathrm{PR}_l\approx l+1$.
\end{proposition}
\begin{proof}
$\E\sum_{ij}(W\T CW)_{ij}^2$ has two contributions. The $n(n-1)$ off-diagonal pairs give $s^4\|C\|_F^2$ each, and the
$n$ diagonal terms give $s^4((\operatorname{tr}C)^2+2\|C\|_F^2)$ each, by Wick. Also $\E\operatorname{tr}W\T CW=s^2n\operatorname{tr}C$.
\end{proof}

Measured along the same network ($n=1024$, $4000$ inputs, finite-sample corrected), $n/\mathrm{PR}(C_{\perp,l})$ is
\begin{center}\small
\begin{tabular}{@{}lccccccccccccccc@{}}\toprule
$l$&1&2&3&4&5&6&7&8&9&10&11&12&13&14&15\\\midrule
$n/\mathrm{PR}$&1.55&2.40&3.25&4.22&5.08&5.93&6.88&7.78&8.41&9.12&9.90&10.76&11.55&12.17&12.91\\\bottomrule
\end{tabular}
\end{center}
The growth is linear with slope $\approx0.85$ instead of $1$, because gating and the higher Mehler terms add diffuse
mass. So $n_{\rm eff}\approx n/(1+0.85\,l)$. The effective expansion parameter of the per-neuron transverse sector is then
$n_{\rm eff}^{-1/2}\approx\sqrt{0.85\,L/n}$: the depth-to-width ratio, not $n^{-1/2}$. The measured skewness RMS fits this
($\sqrt{16/1024}=0.125$ against $0.116$; at $n=512$, $0.177$ against $0.203$). It matches the general finding of
effective-theory analyses of deep networks that finite-width effects grow like $L/n$.

\section{The per-neuron ledger at $n=1024$, $L=16$}
Each entry is $\sigma\,\frac{m_r}{r!}\,\mathrm{He}_{r-2}(-\alpha)\varphi(\alpha)$ (stage~20, Thm.~2.2). The RMS is taken over neurons with
$\alpha\sim N(0,\rho/(1-\rho))$, $\rho=0.916$, $\sigma=0.4$, using the measured standardised cumulants.
\begin{center}\small
\begin{tabular}{@{}lccl@{}}\toprule
term & standardised size & RMS mean shift & verdict\\\midrule
$m_3=\kappa_3$ & $0.116$ & $9.8\times10^{-4}$ & compute to $\approx6\%$\\
$m_4=\kappa_4$, coherent part & $0.066$ & $1.7\times10^{-4}$, same sign & collective resummation\\
$m_4=\kappa_4$, incoherent part & $0.026$ & $6.8\times10^{-5}$ & compute ($K$-register)\\
$m_6\ni10\kappa_3^2$ & $0.135$ & $3.5\times10^{-5}$ & include (free from $\kappa_3$)\\
$m_7\ni35\kappa_3\kappa_4$ (total $\kappa_4$) & $0.268$ & $2.1\times10^{-5}$ & include (free)\\
$m_7\ni35\kappa_3\kappa_4^{\rm inc}$ & $0.106$ & $8\times10^{-6}$ & negligible\\
$m_8\ni35\kappa_4^2$, $m_9\ni280\kappa_3^3$ & --- & $\le3.5\times10^{-6}$ & negligible\\\bottomrule
\end{tabular}
\end{center}
The target RMS is $6\times10^{-5}$. The products of cumulants enter with known Edgeworth coefficients:
$m_6=\kappa_6+10\kappa_3^2$, $m_7=\kappa_7+35\kappa_3\kappa_4$, $m_8=\kappa_8+56\kappa_3\kappa_5+35\kappa_4^2$, all read from the generating function
$\E e^{tF-t^2/2}=\exp\sum_{r\ge3}\kappa_rt^r/r!$. They need no new information beyond $\kappa_3$ and $\kappa_4$. The pure higher
cumulants $\kappa_5,\kappa_6,\kappa_7$ were not measured. The effective-dimension scaling suggests
$\kappa_5\sim n_{\rm eff}^{-3/2}\approx10^{-3}$, which would give a shift of about $10^{-5}$.

\noindent\textbf{Design constraints this imposes.}
\begin{itemize}[leftmargin=1.2em,itemsep=1pt,topsep=1pt]
\item The per-neuron $\kappa_3$ readout is the dominant term and must be accurate to about $6\%$ RMS. By
stage~21's localisation table, $\eta\approx0.06$ needs the last $\ell_3\approx11$ layers of history.
\item The per-neuron readout should be the Edgeworth series in standardised cumulants through $m_7$,
with products included, in the collective form. The collective mixture removes the coherent $\kappa_4$.
\item The incoherent $\kappa_4$ is needed to about $50\%$, which is the incoherent part of $K^z_{kk}$.
\end{itemize}

\section{Adaptive consequences}\label{sec:adapt}
\begin{enumerate}[label=(\roman*),leftmargin=1.6em,itemsep=1pt,topsep=1pt]
\item \emph{Order selector.} Pass~1 computes $n_{\rm eff}=(\operatorname{tr}C_\perp)^2/\|C_\perp\|_F^2$. The scale of the transverse orders is
$n_{\rm eff}^{-1/2}$. Pass~2's probe block of $N=n$ inputs gives $N(N-1)/2$ pair overlaps, hence $\E\rho_\perp^{2k}$ with negligible
error. This checks $n_{\rm eff}$ and supplies the tail term of Theorem~\ref{thm:analytic}.
\item \emph{Skewness calibration test.} The same block, read one layer further, gives per-neuron sample
skewnesses $\hat\gamma_k$ with noise $\approx\sqrt{6/N}=0.077$. Regressing $\hat\gamma_k$ on the theory's
$\kappa_{3,k}/\sigma_k^3$ across the $n$ neurons, with signal RMS $0.116$, estimates the slope with standard error
$\approx0.077/(0.116\sqrt n)\approx0.02$. A slope away from $1$ flags a systematic error in the history readout,
the dominant term, at the $2\%$ level. That means extending the window or adding missing births. Shared
inputs correlate the errors across neurons only weakly, at $O(n_{\rm eff}^{-1/2})$.
\item \emph{Collective check.} The block also gives $\mathrm{Var}\log R^2$ and $\kappa_3(\log R^2)$ (stage~19). The value measured
here, $0.029$, agrees with the carried-vector prediction scale.
\end{enumerate}

\section*{Status}
\noindent\textbf{Exact:} Theorem~\ref{thm:mmd} (with (a) and (b)), Lemma~\ref{lem:pos},
Theorem~\ref{thm:analytic} and Proposition~\ref{prop:pr}.

\noindent\textbf{Measured} (one realised network, quick check): the table of \S\ref{sec:real}.

\noindent\textbf{Estimates:} the per-neuron ledger. These are exact Edgeworth coefficients times measured RMS
sizes, assuming $\alpha$ is independent of the cumulant sizes.

\noindent\textbf{Corrected:} the standardised-size reading of stages~20 and 22 (depth factor
$(1-\rho)^{-(r-2)/2}$).

\noindent\textbf{Open:}
\begin{itemize}[leftmargin=1.2em,itemsep=0pt,topsep=1pt]
\item the slope $0.85$ of $n/\mathrm{PR}$ from the gated Mehler transport, beyond the linear $1$ of
Proposition~\ref{prop:pr};
\item the accuracy of the history $\kappa_3$ at the $6\%$ level;
\item level-wise mismatch bounds $\|\Delta_{jm}\|$ for the certificate.
\end{itemize}

\end{document}
